% !PRINTED 224
\weberpage{224}
%<<CONT
$n$, then one takes two roots $x_{\mu,\mu'}$, $x_{\nu,\nu'}$
into the same row or into different rows, according as

\[
\mu\nu' - \nu\mu' \equiv 0 \ (\mathrm{mod}\ p)
\]

or not; if accordingly $x_{\nu,\nu'}$ and $x_{\nu_1,\nu_1'}$ belong to one row with
$x_{\mu,\mu'}$, then they belong to the same row with one another as well. The
number of the rows so formed is, as is easily shown,
$p + 1$, and in this way one obtains, as the first resolvent of the
division equation, a transformation equation belonging to the divisor $p$.
We shall show later, by another route, how the
transformation equations for composite divisors can be reduced to
those for prime divisors.

\begin{center}
\rule{3cm}{0.4pt}
\end{center}
